% !TeX root = ../thesis.tex
% ===================== ЗАКЛЮЧЕНИЕ =====================
\chapter*{Заключение}
\addcontentsline{toc}{chapter}{Заключение}

В работе разработан и верифицирован программный солвер для прямого решения кинетического уравнения Больцмана
консервативным проекционным методом Черемисина с вычислением интеграла столкновений на 8-мерных сетках Коробова.
Основные результаты:
\begin{enumerate}
  \item реализована общая численная схема проекционного метода; подтверждено машинно-точное сохранение
        массы, импульса и энергии, положительность функции распределения и невозрастание $H$-функции;
  \item на задаче релаксации квази-шока воспроизведены сохранение температуры, релаксация продольной
        компоненты тензора температуры и установление равновесного распределения Максвелла; показана
        многомодовая структура релаксации в сравнении с моделью БГК;
  \item на задаче релаксации бинарной смеси воспроизведены выравнивание парциальных температур к
        равновесию $T=1{,}5$, сохранение концентраций и монотонный рост суммарной энтропии;
  \item решена неоднородная задача переноса о поршне: получена установившаяся ударная волна, скачки
        плотности и температуры ($n_1/n_0=2{,}95$, $T_1/T_0=3{,}62$) и скорость фронта ($V=-1{,}22$)
        согласуются с теорией Ренкина--Гюгонио (невязка $\sim2\,\%$) при машинно-точном сохранении
        инвариантов;
  \item расчёты выполнены на суперкомпьютерном кластере; показана структурная параллелизуемость метода.
\end{enumerate}

\textbf{План дальнейших исследований.}
Основное направление развития~--- промышленная компилируемая реализация солвера с распараллеливанием средствами MPI и
GPU.
Перспективны также обобщение на смеси с различными массами компонентов и переход от модели твёрдых сфер к более
реалистичным потенциалам межмолекулярного взаимодействия.

